\documentclass[french,12pt]{article}

% ============================================================================================
% ======= Configurations 
% ============================================================================================

% ======= Ajout de paquets
% ============================================================================================
% ======= Paquets
% ============================================================================================
\RequirePackage{fancyhdr}
\RequirePackage{fancyvrb}

% ============================================================================================
% ======= Pour le paquet ProfCollege
% ============================================================================================
\usepackage{ProfCollege}

% Pour la gestion des fontes maths notamment en mode LuaLaTeX.
\usepackage{unicode-math}
% Par exemple, une fonte sans serif pour les briques Scratch.
\newfontfamily\myfontScratch[]{FreeSans}

% ============================================================================================
% ======= Le reste
% ============================================================================================

% Pour générer du faux texte
\usepackage{lipsum}

% Pour les hyperliens, notamment des sommaires cliquables
\usepackage{hyperref}


% ======= Ajout de commandes
% ============================================================================================
% ======= 1ere de couverture
% ============================================================================================

\newcommand\myAuthorName{Mireille GAIN et Sébastien LOZANO}
\newcommand\myAuthorSchoolName{Établissements}
\newcounter{postCurrentSchoolYear}
\setcounter{postCurrentSchoolYear}{\the\year}
\addtocounter{postCurrentSchoolYear}{1}
\newcommand\currentSchoolYear{Année  \the\year ~- \thepostCurrentSchoolYear}

% ============================================================================================
% ======= Sommaire
% ============================================================================================
\renewcommand{\contentsname}{Sommaire}           % Changer le nom de la table des matières

% ============================================================================================
% ======= Styles
% ============================================================================================

\fancypagestyle{firstCover}{%   1ere de couverture
    \fancyhf{}%                 On initialise headers and footers à rien du tout !        
    \renewcommand{\headrulewidth}{0pt}%trait horizontal pour l'en-tête
    %\renewcommand{\footrulewidth}{0.4pt}%trait horizontal pour le pied de page    
}
\fancypagestyle{backCover}{%   4ere de couverture
    \fancyhf{}%                 On initialise headers and footers à rien du tout !        
    %\renewcommand{\headrulewidth}{0pt}%trait horizontal pour l'en-tête
    %\renewcommand{\footrulewidth}{0.4pt}%trait horizontal pour le pied de page
}

% ============================================================================================
% Factorisation de commandes
% ============================================================================================


% ======= Pour inclure de façon sélective
\includeonly{%
	./numerationEgyptienne/groupe1.tex,
	./numerationEgyptienne/groupe2.tex,
	./numerationRomaine/groupe1.tex,
	./numerationRomaine/groupe2.tex,
	./numerationShadok/groupe1.tex,
	./numerationShadok/groupe2.tex,
	./numerationInca/groupe1.tex,
	./numerationInca/groupe2.tex,
	./numerationChinoise/groupe1.tex,
	./numerationChinoise/groupe2.tex,
	./numerationBabylonienne/groupe1.tex,
	./numerationBabylonienne/groupe2.tex,
	./numerationBinaire/groupe1.tex,
	./numerationBinaire/groupe2.tex,
	./numerationTriozon/groupe1.tex,
	./numerationTriozon/groupe2.tex,
	./numerationDni/groupe1.tex,
	./numerationDni/groupe2.tex,
}
% ============================================================================================
% ======= Début du document
% ============================================================================================
\begin{document}
    \pagestyle{firstCover}
\begin{titlepage}
    \parindent=0pt
    Texte\hspace*{\stretch{1}} TP numérations \LaTeX 

    Texte\hspace*{\stretch{1}} Paquet ProfCollege
    \vspace*{\stretch{1}}
    \begin{center}
        \includegraphics[scale=0.5]{images/8.png}%
    \end{center}
    \vspace*{\stretch{1}}
    \hrulefill
    \begin{center}\bfseries\Huge
        TP numérations
        
        de Mireille GAIN avec \LaTeX
    \end{center}
    \hrulefill
    \vspace*{1cm}
    \begin{center}\bfseries\Large
        \myAuthorName
    \end{center}
    \begin{center}\bfseries
        \myAuthorSchoolName

        \currentSchoolYear
    \end{center}    
        
    \vspace*{\stretch{2}}
    \begin{flushright}
           Dernière mise à jour le \today 
    \end{flushright}   
    \newcommand\myBubbleColor{black}
    \begin{tikzpicture}[remember picture, overlay]
     \begin{scope}[shift={(current page.south west)},shift={(-.5,-.5)},scale=1]
     \shade[ball color=\myBubbleColor,opacity=.6] (0,0) circle (10ex);
     \shade[ball color=\myBubbleColor,opacity=.8] (1.7,1) circle (5ex);
     \shade[ball color=\myBubbleColor,opacity=.8] (1.5,3) circle (2ex);
     \shade[ball color=\myBubbleColor,opacity=.5] (-0.5,3) circle (1ex);
     \shade[ball color=\myBubbleColor,opacity=.8] (1,4) circle (1ex);
     \shade[ball color=\myBubbleColor,opacity=.6] (3.5,2.5) circle (2ex);
     \shade[ball color=\myBubbleColor,opacity=.8] (2.5,4.5) circle (4ex);
     \shade[ball color=\myBubbleColor,opacity=.5] (3,4) circle (3ex);
     \shade[ball color=\myBubbleColor,opacity=.8] (4.5,4.5) circle (3ex);
     \shade[ball color=\myBubbleColor,opacity=.5] (5.1,4.7) circle (2ex);
     \shade[ball color=\myBubbleColor,opacity=.8] (5,6) circle (1.5ex);
     \shade[ball color=\myBubbleColor,opacity=.6] (3.5,5.5) circle (2ex);
     \shade[ball color=\myBubbleColor,opacity=.8] (5,3) circle (1ex);
     \end{scope}
     \end{tikzpicture}
\end{titlepage}            				% On inclut la 1ere de couverture

    \tableofcontents									% Sommaire
	
% ======= On inclut la numération Egyptienne
	\def\currentpath{./numerationEgyptienne}            % On définit le répertoire courant
	\clearpage
	\section{Numération Babylonienne}
Début commun du doc
\lipsum[1-2]	
	\subsection{Spécifique Groupe1}
	\subsection{Spécifique Groupe2}
	fin commune du doc

% ======= On inclut la numération Romaine
	\def\currentpath{./numerationRomaine}            	% On définit le répertoire courant
	\clearpage
	\section{Numération Babylonienne}
Début commun du doc
\lipsum[1-2]	
	\subsection{Spécifique Groupe1}
	\subsection{Spécifique Groupe2}
	fin commune du doc

% ======= On inclut la numération Shadok
	\def\currentpath{./numerationShadok}            	% On définit le répertoire courant
	\clearpage
	\section{Numération Babylonienne}
Début commun du doc
\lipsum[1-2]	
	\subsection{Spécifique Groupe1}
	\subsection{Spécifique Groupe2}
	fin commune du doc

% ======= On inclut la numération Inca
	\def\currentpath{./numerationInca}            		% On définit le répertoire courant
	\clearpage
	\section{Numération Babylonienne}
Début commun du doc
\lipsum[1-2]	
	\subsection{Spécifique Groupe1}
	\subsection{Spécifique Groupe2}
	fin commune du doc

% ======= On inclut la numération Chinoise
	\def\currentpath{./numerationChinoise}            	% On définit le répertoire courant
	\clearpage
	\section{Numération Babylonienne}
Début commun du doc
\lipsum[1-2]	
	\subsection{Spécifique Groupe1}
	\subsection{Spécifique Groupe2}
	fin commune du doc

% ======= On inclut la numération Babylonienne
	\def\currentpath{./numerationBabylonienne}           % On définit le répertoire courant
	\clearpage
	\section{Numération Babylonienne}
Début commun du doc
\lipsum[1-2]	
	\subsection{Spécifique Groupe1}
	\subsection{Spécifique Groupe2}
	fin commune du doc

% ======= On inclut la numération Binaire
	\def\currentpath{./numerationBinaire}            	% On définit le répertoire courant
	\clearpage
	\section{Numération Babylonienne}
Début commun du doc
\lipsum[1-2]	
	\subsection{Spécifique Groupe1}
	\subsection{Spécifique Groupe2}
	fin commune du doc

% ======= On inclut la numération Triozon
	\def\currentpath{./numerationTriozon}            	% On définit le répertoire courant
	\clearpage
	\section{Numération Babylonienne}
Début commun du doc
\lipsum[1-2]	
	\subsection{Spécifique Groupe1}
	\subsection{Spécifique Groupe2}
	fin commune du doc

% ======= On inclut la numération D'ni
	\def\currentpath{./numerationDni}            		% On définit le répertoire courant
	\clearpage
	\section{Numération Babylonienne}
Début commun du doc
\lipsum[1-2]	
	\subsection{Spécifique Groupe1}
	\subsection{Spécifique Groupe2}
	fin commune du doc

    \pagestyle{backCover}
\parindent=0pt
Texte \hspace*{\stretch{1}} TP numérations \LaTeX 

Texte \hspace*{\stretch{1}} Paquet ProfCollege
\vspace*{\stretch{1}}
\begin{center}
    \includegraphics[scale=0.5]{images/8.png}%
\end{center}
\vspace*{\stretch{1}}
\hrulefill

Du texte ici si je veux \ldots

\hrulefill
\vspace*{1cm}
\begin{center}\bfseries\Large
    \myAuthorName
\end{center}
\begin{center}\bfseries
    \myAuthorSchoolName

    \currentSchoolYear
\end{center}    
    
\vspace*{\stretch{2}}
\begin{flushright}
       Dernière mise à jour le \today 
\end{flushright}   
          				% On inclut la 4e de couverture
\end{document}